% Part: first-order-logic
% Chapter: beyond
% Section: second-order-logic

\documentclass[../../../include/open-logic-section]{subfiles}

\begin{document}

\olfileid{fol}{byd}{sol}

\olsection{দ্বিতীয়-ক্রমের যুক্তিবিদ্যা}

দ্বিতীয়-ক্রমের যুক্তিবিদ্যার ভাষায় শুধু ব্যক্তিবস্তুর কোনো
!!a{domain}-এর উপর নয়, সেই !!{domain}-এর সম্পর্কগুলির উপরও পরিমাণায়ন
করা যায়। কোনো প্রথম-ক্রমের ভাষা $\Lang{L}$ দেওয়া থাকলে প্রত্যেক
$k$-এর জন্য এমন !!{variable}s $R$ যোগ করা হয় যেগুলি $k$-স্থানীয়
সম্পর্কের উপর বিস্তৃত, এবং ওই চলরাশিগুলির উপর পরিমাণায়নের অনুমতি দেওয়া
হয়। $R$ যদি কোনো $k$-স্থানীয় সম্পর্কের !!a{variable} হয় এবং
$t_1$, \dots,~$t_k$ যদি সাধারণ (প্রথম-ক্রমীয়) পদ হয়, তবে
$\Atom{R}{t_1,\dots,t_k}$ একটি পরমাণু !!{formula}। অন্য সব ক্ষেত্রে
!!{formula}s-এর সেট প্রথম-ক্রমের যুক্তিবিদ্যার মতোই সংজ্ঞায়িত হয়,
শুধু দ্বিতীয়-ক্রমের পরিমাণায়নের জন্য অতিরিক্ত বিধি থাকে। লক্ষ করো,
শুধু প্রথম-ক্রমীয় পদের জন্যই আমাদের !!{identity} আছে: $R$ ও $S$ একই
স্থানসংখ্যা~$k$-এর সম্পর্ক-!!{variable}s হলে $\eq[R][S]$-কে নিচেরটির
সংক্ষিপ্ত রূপ হিসেবে সংজ্ঞায়িত করা যায়:
\[
\lforall[x_1 \dots][\lforall[x_k][(\Atom{R}{x_1, \dots, x_k} \liff
  \Atom{S}{x_1, \dots, x_k})]].
\]

দ্বিতীয়-ক্রমের যুক্তিবিদ্যার বিধিগুলি নতুন দ্বিতীয়-ক্রমীয় চলরাশির
ক্ষেত্রে পরিমাণসূচকের বিধিগুলিকে শুধু প্রসারিত করে। তবে এই চলরাশিগুলি
$\Lang{L}$-এর !!{predicate}s এবং আরও সাধারণভাবে $\Lang{L}$-এর
!!{formula}s-এর সঙ্গে কীভাবে মিথস্ক্রিয়া করে, তা একটু সতর্কভাবে
ব্যাখ্যা করতে হয়। ন্যূনতমভাবে সম্পর্ক-চলরাশিগুলি পদ হিসেবে গণ্য হয়,
তাই নিচের রূপের অনুমান থাকে:
\[
!A(R) \Proves \lexists[R][!A(R)]
\]
কিন্তু $\Lang{L}$ যদি পাটীগণিতের ভাষা হয় এবং তাতে ধ্রুব সম্পর্ক-সংকেত
$<$ থাকে, তবে নিচের অনুমানটিও বৈধ হবে বলে আশা করা যায়:
\[
x < y \Proves \lexists[R][\Atom{R}{x,y}]
\]
অথবা কোনো দেওয়া !!{formula}~$!A$-এর জন্য,
\[
!A(x_1, \dots, x_k) \Proves \lexists[R][\Atom{R}{x_1,\dots,x_k}]
\]
আরও সাধারণভাবে নিচের রূপের অনুমান গ্রহণ করতে চাইতে পারি:
\[
\Subst{!A}{\lambd[\vec x][!B(\vec x)]}{R} \Proves
\lexists[R][!A]
\]
এখানে $\Subst{!A}{\lambd[\vec x][!B(\vec x)]}{R}$ দিয়ে~$!A$-তে
$\Obj{R}{t_1,\dots,t_k}$ রূপের প্রত্যেক পরমাণু !!{formula}-কে
$!B(t_1, \dots, t_k)$ দিয়ে প্রতিস্থাপনের ফল বোঝায়। শেষ বিধিটি একটি
{\em ধর্মনির্দেশ-ছক} থাকার সমতুল্য; অর্থাৎ, নিচের রূপের একটি করে
স্বতঃসিদ্ধ
\[
\lexists[R][\lforall[x_1, \dots, x_k][(!A(x_1, \dots, x_k) \liff
\Atom{R}{x_1, \dots, x_k})]],
\]
দ্বিতীয়-ক্রমের ভাষার প্রত্যেক !!{formula}~$!A$-এর জন্য, যেখানে $R$
মুক্ত চলরাশি নয়। (অনুশীলনী: দেখাও, $!A$-তে $R$ ঘটার অনুমতি দিলে এই
ছকটি অসঙ্গত!)

যুক্তিবিদেরা “দ্বিতীয়-ক্রমের যুক্তিবিদ্যার স্বতঃসিদ্ধ” বলতে সাধারণত
দ্বিতীয়-ক্রমীয় পরিমাণসূচক-বিধি ও ধর্মনির্দেশ-ছক দিয়ে প্রথম-ক্রমের
যুক্তিবিদ্যার ন্যূনতম সম্প্রসারণটি বোঝান। তবে এই স্বতঃসিদ্ধ ও বিধিগুলির
দুর্বলতর উপব্যবস্থাও অধ্যয়ন করা প্রায়ই আকর্ষণীয়। যেমন, পূর্ণ
সাধারণতায় ধর্মনির্দেশের স্বতঃসিদ্ধ-ছকটি \emph{অপ্রেডিকেটিভ}: এটি
দ্বিতীয়-ক্রমীয় পরিমাণসূচকযুক্ত !!a{formula} দ্বারা “সংজ্ঞায়িত”
$\Atom{R}{x_1, \dots, x_k}$ সম্পর্কের অস্তিত্ব দাবি করতে দেয়; আর এই
পরিমাণসূচকগুলি এমন সব সম্পর্কের সেটের উপর বিস্তৃত---যে সেটে $R$
নিজেও আছে!{} বিংশ শতাব্দীর শুরুর দিকে রাসেলের কূটাভাসের একটি প্রচলিত
প্রতিক্রিয়া ছিল এমন সংজ্ঞাকেই দোষ দেওয়া এবং গণিতের ভিত্তি নির্মাণে
সেগুলি এড়ানো। !!{formula}~$!A$-তে দ্বিতীয়-ক্রমীয় পরিমাণসূচকের ব্যবহার
নিষিদ্ধ করলে ধর্মনির্দেশের একটি {\em প্রেডিকেটিভ} রূপ পাওয়া যায়, যা
কিছুটা দুর্বলতর।

অর্থতাত্ত্বিক দৃষ্টিতে কোনো দ্বিতীয়-ক্রমীয় !!{structure}-কে ভাষাটির
একটি প্রথম-ক্রমীয় !!{structure} এবং !!{domain}-এর উপর এমন সম্পর্কের
একটি সেট নিয়ে গঠিত ভাবা যায় যার উপর দ্বিতীয়-ক্রমীয় পরিমাণসূচকগুলি
বিস্তৃত (আরও নির্ভুলভাবে, প্রত্যেক $k$-এর জন্য $k$ স্থানসংখ্যার
সম্পর্কগুলির একটি সেট থাকে)। অবশ্য ধর্মনির্দেশকে !!{derivation}
ব্যবস্থায় অন্তর্ভুক্ত করলে আরও দাবি করতে হয় যে “দ্বিতীয়-ক্রমীয়
অংশে” ধর্মনির্দেশের স্বতঃসিদ্ধগুলি পরিতৃপ্ত করার মতো যথেষ্ট সম্পর্ক
আছে---তা না হলে !!{derivation} ব্যবস্থাটি নির্ভুল নয়!{} যথেষ্ট সম্পর্ক
নিশ্চিত করার একটি সহজ উপায় হলো দ্বিতীয়-ক্রমীয় অংশে প্রথম-ক্রমীয়
অংশের উপর \emph{সব} সম্পর্ক নেওয়া। এমন !!a{structure}-কে
\emph{পূর্ণ} বলা হয় এবং এক অর্থে এটিই ভাষাটির “অভিপ্রেত
!!{structure}”। শুধু পূর্ণ !!{structure}s-এ মনোযোগ সীমিত করলে
\emph{পূর্ণ} দ্বিতীয়-ক্রমীয় অর্থতত্ত্ব পাওয়া যায়। তখন কোনো
!!{structure} নির্দিষ্ট করা মূলত তার প্রথম-ক্রমীয় অংশটি নির্দিষ্ট
করাতেই নেমে আসে, কারণ দ্বিতীয়-ক্রমীয় অংশের বিষয়বস্তু তা থেকে
অন্তর্নিহিতভাবে নির্ধারিত হয়।

সংক্ষেপে, দ্বিতীয়-ক্রমের যুক্তিবিদ্যা নিয়ে কথায় কিছু দ্ব্যর্থতা
আছে। !!{derivation} ব্যবস্থার দিক থেকে নিচের যেকোনোটি মনে থাকতে পারে:
\begin{enumerate}
\item কয়েকটি ধর্মনির্দেশ-স্বতঃসিদ্ধসহ একটি “ন্যূনতম” দ্বিতীয়-ক্রমীয়
  !!{derivation} ব্যবস্থা।
\item পূর্ণ ধর্মনির্দেশসহ “মানক” দ্বিতীয়-ক্রমীয় !!{derivation}
  ব্যবস্থা।
\end{enumerate}
অর্থতত্ত্বের দিক থেকেও নিচের যেকোনোটিতে আগ্রহ থাকতে পারে:
\begin{enumerate}
\item “দুর্বল” অর্থতত্ত্ব, যেখানে প্রথম-ক্রমীয় একটি অংশ এবং
  ধর্মনির্দেশের স্বতঃসিদ্ধগুলি পরিতৃপ্ত করার মতো যথেষ্ট বড় একটি
  দ্বিতীয়-ক্রমীয় অংশ নিয়ে !!a{structure} গঠিত।
\item “মানক” দ্বিতীয়-ক্রমীয় অর্থতত্ত্ব, যেখানে শুধু পূর্ণ
  !!{structure}s বিবেচনা করা হয়।
\end{enumerate}
কোন !!{derivation} ব্যবস্থা বা অর্থতত্ত্ব তাঁদের মনে আছে তা
যুক্তিবিদেরা নির্দিষ্ট না করলে সাধারণত দুটি তালিকারই দ্বিতীয় দফাটি
বোঝান। এই অর্থতত্ত্ব ব্যবহারের সুবিধা হলো, পরে আমরা যেমন দেখব, এটি
অনেক স্বাভাবিক গাণিতিক গঠনের সমরূপতা-অবধি একক বর্ণনা দেয়; একই সঙ্গে
!!{derivation} ব্যবস্থাটি বেশ শক্তিশালী এবং এই অর্থতত্ত্বের জন্য
নির্ভুল। অসুবিধা হলো, !!{derivation} ব্যবস্থাটি অর্থতত্ত্বটির জন্য
\emph{সম্পূর্ণ নয়}; বস্তুত পূর্ণ দ্বিতীয়-ক্রমীয় অর্থতত্ত্বের জন্য
\emph{কোনো} কার্যকরভাবে দেওয়া !!{derivation} ব্যবস্থাই সম্পূর্ণ নয়।
অন্যদিকে পরে আমরা দেখব, দুর্বলতর অর্থতত্ত্বের জন্য !!{derivation}
ব্যবস্থাটি \emph{সম্পূর্ণ}; ফলে কোনো বাক্য প্রমাণযোগ্য না হলে এমন
\emph{কোনো} !!{structure} আছে---পূর্ণটি হওয়া আবশ্যক নয়---যেখানে
বাক্যটি মিথ্যা।

দ্বিতীয়-ক্রমের যুক্তিবিদ্যার ভাষাটি খুবই সমৃদ্ধ। এক-স্থানীয় সম্পর্ককে
!!{domain}-এর উপসেটের সঙ্গে সনাক্ত করা যায়, তাই বিশেষ করে এই সেটগুলির
উপর পরিমাণায়ন করা যায়; যেমন, একটি মাত্র স্বতঃসিদ্ধ দিয়ে স্বাভাবিক
সংখ্যার জন্য আরোহ প্রকাশ করা যায়:
\[
\lforall[R][((\Atom{R}{\Obj{0}} \land \lforall[x][(\Atom{R}{x} \lif
    \Atom{R}{x'})]) \lif \lforall[x][\Atom{R}{x}])].
\]
পাটীগণিতের ভাষায় $\Obj 0, \Obj \prime, +, \times$ ও $<$ সংকেতগুলি
আছে ধরলে তাদের আচরণ বর্ণনা করতে নিচের স্বতঃসিদ্ধগুলি যোগ করা যায়:
\begin{enumerate}
\item $\lforall[x][\lnot x' = \Obj 0]$
\item $\lforall[x][\lforall[y][(x' = y' \lif x = y)]]$
\item $\lforall[x][(x + \Obj 0) = x]$
\item $\lforall[x][\lforall[y][(x + y') = (x + y)']]$
\item $\lforall[x][(x \times \Obj 0) = \Obj 0]$
\item $\lforall[x][\lforall[y][(x \times y') = ((x \times y) + x)]]$
\item $\lforall[x][\lforall[y][(x < y \liff \lexists[z][y = (x + z')])]]$
\end{enumerate}
% BN-SRC-110 TARGET-CORRECTION-BEGIN
\par\smallskip\noindent\textnormal{\small\textbf{উৎস-সংশোধন (BN-SRC-110):} হিমায়িত ইংরেজি উৎসে ভাষাটির ঘোষিত উত্তরসূরি সংকেত $\Obj\prime$ ও অন্য সব স্বতঃসিদ্ধের প্রাইম-লিপির বদলে দ্বিতীয় স্বতঃসিদ্ধে অঘোষিত $s(x)$ ও $s(y)$ লেখা ছিল। বাংলা সূত্রে একই ঘোষিত লিপি অনুসারে $x'$ ও $y'$ পুনরুদ্ধার করা হয়েছে।} % BN-SRC-110 TARGET-CORRECTION-END
পূর্ণ দ্বিতীয়-ক্রমীয় অর্থতত্ত্ব ব্যবহার করলে উপরের আরোহ-স্বতঃসিদ্ধসহ
এই স্বতঃসিদ্ধগুলি পাটীগণিতের মানক মডেল !!{structure}~$\Struct{N}$-এর
সমরূপতা-অবধি একক বর্ণনা দেয়, তা দেখানো কঠিন নয়। এই স্বতঃসিদ্ধগুলি
সত্য এমন যেকোনো !!{structure}~$\Struct{M}$ দেওয়া থাকলে $\Nat$-এর উপর
সাধারণ পুনরাবৃত্তি দিয়ে $\Nat$ থেকে $\Struct{M}$-এর !!{domain}-এ
একটি অপেক্ষক~$f$ সংজ্ঞায়িত করো, যেন $f(0) = \Assign{\Obj 0}{M}$ এবং
$f(x+1) = \Assign{\prime}{M}(f(x))$। $\Nat$-এর উপর সাধারণ আরোহ এবং
$\Struct M$-এ স্বতঃসিদ্ধ (1) ও~(2) সত্য---এই তথ্য ব্যবহার করে দেখা যায়
$f$ !!{injective}। $f$ যে !!{surjective} তা দেখতে~$P$ হিসেবে
$\Domain{M}$-এর সেই উপাদানগুলির সেট নাও যেগুলি~$f$-এর মানপরিসরে আছে।
$\Struct M$ পূর্ণ বলে $P$ দ্বিতীয়-ক্রমীয় !!{domain}-এ আছে। $f$-এর
নির্মাণ থেকে জানি, $\Assign{\Obj 0}{M}$, $P$-তে আছে এবং
$\Assign{\prime}{M}$-এর অধীনে $P$ বদ্ধ। $\Struct M$-এ
আরোহ-স্বতঃসিদ্ধ সত্য (বিশেষত~$P$-এর জন্য) বলে $P$, $\Struct M$-এর
সমগ্র প্রথম-ক্রমীয় !!{domain}-এর সমান। এতে প্রমাণ হয়, $f$
!!a{bijection}। $\Nat$-এর উপর সাধারণ আরোহ বারবার ব্যবহার করে $f$ যে
একটি হোমোমরফিজম, তা দেখানোও এর চেয়ে কঠিন নয়।

সেটতাত্ত্বিক ভাষায় অপেক্ষক কেবল বিশেষ ধরনের সম্পর্ক; যেমন, এক-স্থানীয়
অপেক্ষক~$f$-কে $\lforall[x][\lexists![y][R(x,y)]]$ পরিতৃপ্তকারী
দ্বি-স্থানীয় সম্পর্ক~$R$-এর সঙ্গে সনাক্ত করা যায়। ফলে অপেক্ষকের উপরও
পরিমাণায়ন করা যায়। তখন পূর্ণ অর্থতত্ত্ব ব্যবহার করে অসীম
!!{structure}s-এর শ্রেণিকে সেইসব !!{structure}s $\Struct M$-এর শ্রেণি
হিসেবে সংজ্ঞায়িত করা যায় যাদের $\Struct M$-এর !!{domain} থেকে তারই একটি যথার্থ
উপসেটে কোনো এক-এক অপেক্ষক আছে:
\[
\lexists[f][(\lforall[x][\lforall[y][(\eq[f(x)][f(y)] \lif \eq[x][y])]]
  \land \lexists[y][\lforall[x][\eq/[f(x)][y]]])].
\]
তখন এই বাক্যটির নঞর্থক রূপটি সসীম !!{structure}s-এর শ্রেণিকে
সংজ্ঞায়িত করে।

এ ছাড়াও রৈখিক ক্রমের সংজ্ঞায় নিচেরটি যোগ করে সুক্রমগুলির শ্রেণি
সংজ্ঞায়িত করা যায়:
\[
\lforall[P][(\lexists[x][\Atom{P}{x}] \lif \lexists[x][(\Atom{P}{x}
    \land \lforall[y][(y < x \lif \lnot \Atom{P}{y})])])].
\]
“সেট”-কে “এক-স্থানীয় সম্পর্ক”-এর সঙ্গে সনাক্ত করলে এটি দাবি করে,
প্রত্যেক অশূন্য সেটের ক্ষুদ্রতম একটি উপাদান আছে। আরেকটি উদাহরণ হিসেবে
গ্রাফের সংযুক্ততার ধারণাটি প্রকাশ করা যায় এই বলে যে, শীর্ষবিন্দুগুলিকে
অসংযুক্ত অংশে ভাগ করে এমন কোনো অ-তুচ্ছ পৃথকীকরণ নেই:
\[
\lnot \lexists[A][(\lexists[x][A(x)] \land \lexists[y][\lnot A(y)]
  \land \lforall[w][\lforall[z][((\Atom{A}{w} \land \lnot \Atom{A}{z})
      \lif \lnot \Atom{R}{w,z})]])].
\]
আরও একটি উদাহরণ হিসেবে এমন সসীম !!{structure}s-এর শ্রেণি সংজ্ঞায়িত
করার চেষ্টা করতে পারো যাদের !!{domain}-এর আকার জোড়। আরও উল্লেখযোগ্য
বিষয় হলো, মূলদ সংখ্যাগুলিকে ধারণকারী সম্পূর্ণ ক্রমিত ক্ষেত্র হিসেবে
বাস্তব সংখ্যার সমরূপতা-অবধি একক বর্ণনা দেওয়া যায়।

সংক্ষেপে, দ্বিতীয়-ক্রমের যুক্তিবিদ্যা প্রথম-ক্রমের যুক্তিবিদ্যার চেয়ে
অনেক বেশি প্রকাশক্ষম। এটি সুসংবাদ; এবার দুঃসংবাদ। আগেই বলেছি, পূর্ণ
দ্বিতীয়-ক্রমীয় অর্থতত্ত্বের জন্য সম্পূর্ণ কোনো কার্যকর
!!{derivation} ব্যবস্থা নেই। ভালো হোক বা মন্দ, প্রথম-ক্রমের
যুক্তিবিদ্যার অনেক ধর্মই এখানে অনুপস্থিত; এর মধ্যে সংহতি এবং
ল্যোভেনহাইম--স্কোলেম উপপাদ্যগুলিও আছে।

অন্যদিকে পূর্ণ দ্বিতীয়-ক্রমীয় অর্থতত্ত্ব ছেড়ে দুর্বলতরটি গ্রহণ করতে
রাজি হলে ন্যূনতম দ্বিতীয়-ক্রমীয় !!{derivation} ব্যবস্থা এই
অর্থতত্ত্বের জন্য সম্পূর্ণ। অন্যভাবে বললে, $\Proves$-কে “ন্যূনতম
ব্যবস্থায় প্রমাণ করে” এবং $\Entails$-কে “দুর্বলতর অর্থতত্ত্বে যৌক্তিক
ফলশ্রুতি” হিসেবে পড়লে দেখানো যায়: যখনই $\Gamma \Entails !A$, তখন
$\Gamma \Proves !A$। নির্দিষ্ট ধর্মনির্দেশ-স্বতঃসিদ্ধগুলি
!!{derivation} ব্যবস্থায় রাখতে চাইলে অর্থতত্ত্বকে ওই স্বতঃসিদ্ধগুলি
পরিতৃপ্তকারী দ্বিতীয়-ক্রমীয় গঠনগুলিতে সীমিত করতে হবে। যেমন,
$\Delta$ যদি ধর্মনির্দেশ-স্বতঃসিদ্ধগুলির একটি সেট (সম্ভবত সবগুলিই) হয়,
তবে $\Gamma \cup \Delta \Entails !A$ হলে $\Gamma \cup \Delta
\Proves !A$। বিশেষত, বিবেচ্য ধর্মনির্দেশ-স্বতঃসিদ্ধগুলি ব্যবহার করে
$!A$ প্রমাণযোগ্য না হলে $\lnot !A$-এর এমন একটি মডেল আছে যেখানে এই
ধর্মনির্দেশ-স্বতঃসিদ্ধগুলি তবু সত্য।

দুর্বলতর অর্থতত্ত্বের জন্য সম্পূর্ণতা উপপাদ্যটি সত্য দেখার সবচেয়ে সহজ
উপায় হলো দ্বিতীয়-ক্রমের যুক্তিবিদ্যাকে নিচের মতো একটি বহুজাতীয়
যুক্তিবিদ্যা হিসেবে ভাবা। একটি জাতিকে সাধারণ “প্রথম-ক্রমীয়”
!!{domain} হিসেবে ব্যাখ্যা করা হয়; এরপর প্রত্যেক $k$-এর জন্য থাকে
“$k$ স্থানসংখ্যার সম্পর্ক”-এর !!a{domain}। ভাষাটিতে অন্তর্নির্মিত
সম্পর্ক-সংকেত “$\Atom{\Obj{true}_k}{R,x_1,\dots,x_k}$” নেওয়া হয়; এর
উদ্দেশ্য হলো বলা যে $R$, $x_1$, \dots,~$x_k$-এর জন্য সত্য, যেখানে
$R$ হলো “$k$-স্থানীয় সম্পর্ক” জাতির চলরাশি এবং $x_1$, \dots,~$x_k$
হলো প্রথম-ক্রমীয় জাতির বস্তু।

এই সনাক্তকরণে দুর্বল দ্বিতীয়-ক্রমীয় অর্থতত্ত্ব মূলত বহুজাতীয়
যুক্তিবিদ্যার পরিচিত অর্থতত্ত্ব; আর বহুজাতীয় যুক্তিবিদ্যাকে যে
প্রথম-ক্রমের যুক্তিবিদ্যায় নিবেশিত করা যায়, তা আমরা ইতিমধ্যে দেখেছি।
দুই দিকের অনুবাদগুলি ধরলে দ্বিতীয়-ক্রমের যুক্তিবিদ্যার দুর্বল ধারণাটি
তাই আসলে ছদ্মবেশী প্রথম-ক্রমের যুক্তিবিদ্যার একটি রূপ, যেখানে উপযুক্ত
স্বতঃসিদ্ধ দ্বারা পরিচালিত “বস্তু” ও “সম্পর্ক” উভয়ই !!{domain}-এ আছে।

\end{document}
